% Section 9.3, from lectures/L09/README.md: the objectives, the questions,
% the further reading, and what the next chapter does with all of it.


\section{Review}
\label{c9:sec:review}

\needspace{6\baselineskip}
\subsection*{What you should be able to do now}
After this chapter, you should be able to:
\begin{itemize}
  \item Draw a differential pair and say what its tail current fixes.
  \item Derive the differential gain and account for both factors of two in it.
  \item Derive the common-mode gain and say why the tail appears doubled.
  \item State CMRR, and say which components in the circuit do \textbf{not} affect it.
  \item Choose a tail arrangement from a CMRR requirement, and justify the supply voltage it needs.
  \item Give the two independent reasons a current-mirror load raises the gain.
  \item State the pair's linear input range in millivolts, and say what it does outside it.
  \item Compute an input-referred offset from a stated mismatch, and say which mismatch dominates.
\end{itemize}

\needspace{6\baselineskip}
\subsection*{Questions to test yourself}
You should be able to answer:
\begin{enumerate}
  \item A pair with a 2 mA tail: what is $r_e$ on each side, and why is it 26 ohm rather than 13?
  \item The differential gain to one collector is half what it is to both. Where did the other half go, and what would recover it?
  \item Two designers argue about CMRR. One wants a larger collector resistor. Who is right?
  \item What resistance would the tail need for 80 dB of rejection, and what supply voltage does that imply? Now say what is used instead.
  \item A current-mirror load raises the gain by a factor of ten here. Split that ten into its two causes.
  \item An input signal of 50 mV differential. Is the small-signal gain the right answer? Show why not.
  \item A 1 per cent mismatch between the two collector resistors: what offset does it produce at the input, and how does that compare with 1 mV of $V_{BE}$ mismatch?
\end{enumerate}

\needspace{6\baselineskip}
\subsection*{Further reading}
\begin{itemize}
  \item \emph{The Art of Electronics}, Horowitz and Hill, chapter 2, sections on the differential amplifier, for a second treatment of the same material.
\end{itemize}

\needspace{6\baselineskip}
\subsection*{Looking ahead}
\begin{itemize}
  \item Chapter~\ref{c10:ch:amplifier} assembles Chapter~\ref{c6:ch:bias} to this chapter into one operational amplifier, opens the black box of Chapter~\ref{c3:ch:filters}, closes the loop with Chapter~\ref{c4:ch:feedback}'s arithmetic, and is the capstone.
\end{itemize}
